\emph{Setting.} For a vector $x$, $\tilde x = x$, so $\langle (x\psi)^{\sim}\varphi \rangle_{0} = \langle \tilde\psi\,x\,\varphi \rangle_{0}$: left multiplication by $x$ is symmetric for the product, and $g(x)^{2}\psi = x^{2}\psi = Q(x)\psi$. For a unit vector $x$ it is an involution with eigenspaces $\pi(\pm x)S$, where $\pi(x) = \tfrac12(1 + x)$ satisfies $\pi^{2} = \pi = \tilde\pi$ and $x\pi = \pi x = \pi$; a unit vector orthogonal to $x$ anticommutes with $x$ and exchanges the two eigenspaces, so each has dimension two and $G_{x}$ has signature $(2,2)$ for $x \neq 0$. The product is positive definite on $S$ (computed: the four elements $p$, $e_{1}p$, $e_{2}p$, $e_{1}e_{2}p$ are orthogonal of square $\tfrac12$). By the table of \claimref{std-algebra.periodicity}, $\Cl(3,0)$, of dimension eight, is the algebra of $2 \times 2$ matrices over a division algebra of real dimension two, whose least module has real dimension four, the dimension of $S$.

\emph{(1).} Let $v, v'$ lie in the $+1$ eigenspace of $e_{3}$ on $S$, which is spanned by $p$ and $e_{1}e_{2}p$ ($e_{3}p = p$ and $e_{1}e_{2}$ commutes with $e_{3}$). Then
\[
\langle (\pi v)^{\sim}\,\pi v' \rangle_{0} = \langle \tilde v\,\pi\,v' \rangle_{0} = \tfrac12\langle \tilde v v' \rangle_{0} + \tfrac12\sum_{j} x_{j}\langle \tilde v\,e_{j}\,v' \rangle_{0} = \tfrac12(1 + x_{3})\,\langle \tilde v v' \rangle_{0},
\]
because $e_{1}v'$ and $e_{2}v'$ lie in the $-1$ eigenspace of the symmetric map $e_{3}$, orthogonal to $v$. With $\langle \tilde p\,p \rangle_{0} = \langle (e_{1}e_{2}p)^{\sim}e_{1}e_{2}p \rangle_{0} = \tfrac12$ and $\langle \tilde p\,e_{1}e_{2}\,p \rangle_{0} = 0$ this gives the squares $\tfrac14(1 + x_{3})$ and the orthogonality of $n_{1}, n_{2}$; they lie in $\pi S = G_{+}$, which has rank two, so they are a frame where $x_{3} \neq -1$. The same computation in the $-1$ eigenspace of $e_{3}$, spanned by $e_{1}p$ and $e_{2}p$, gives $s_{1}, s_{2}$ with squares $\tfrac14(1 - x_{3})$. On the equator, $\pi = \pi x$ and $x = \cos a\,e_{1} + \sin a\,e_{2}$, so
\[
\pi e_{1}p = \pi\,x\,e_{1}p = \pi(\cos a + \sin a\,e_{2}e_{1})p = \cos a\,\pi p - \sin a\,\pi e_{1}e_{2}p, \qquad
\pi e_{2}p = \pi\,x\,e_{2}p = \pi(\cos a\,e_{1}e_{2} + \sin a)p = \sin a\,\pi p + \cos a\,\pi e_{1}e_{2}p .
\]

\emph{(2).} $R_{a} = \cos\tfrac a2 + \sin\tfrac a2 f_{1}f_{2} = f_{1}\bigl(\cos\tfrac a2\,f_{1} + \sin\tfrac a2\,f_{2}\bigr)$ is a product of two unit vectors, a rotor in the sense of (L). With this $R_{a}$, $\tilde R_{a} = \cos\tfrac a2 - \sin\tfrac a2 f_{1}f_{2}$, $R_{a}\tilde R_{a} = 1$, and $f_{1}f_{2}$ anticommutes with $f_{1}$ and $f_{2}$ and commutes with every vector orthogonal to both. So $R_{a}f_{1}\tilde R_{a} = R_{a}^{2}f_{1} = (\cos a + \sin a\,f_{1}f_{2})f_{1} = \cos a\,f_{1} - \sin a\,f_{2}$, likewise $R_{a}f_{2}\tilde R_{a} = \sin a\,f_{1} + \cos a\,f_{2}$, and $R_{a}$ fixes the other directions: this is the turn of (1), on the equator $n_{j}$ and $s_{j}$ having the same length. $R_{0} = 1$, $R_{2\pi} = \cos\pi = -1$. On $S^{(k)}$ the frames are $k$ copies of those of (1), mutually orthogonal, the clutching function is the same turn in each of the $k$ planes, and the product of the $k$ commuting rotors $\cos\tfrac a2 + \sin\tfrac a2 f_{2j-1}f_{2j}$ induces it and ends at $(-1)^{k}$.

\emph{(L), the rotors.} In $\Cl(n,0)$ a unit vector has $Q(u) = 1$, so an even versor of unit vectors $R = u_{1}\cdots u_{2j}$ has $R\tilde R = \prod Q(u_{i}) = 1$: $\Spin(n,0) = \Spin^{+}(n,0)$ in the notation of Definition \claimref{obj.std-algebra.pin-spin}, and $\mathrm{SO}^{+}(n,0)$, the products of evenly many reflections with $\prod Q(u_{i}) > 0$, is all of $\mathrm{SO}(n)$. On an even element the twisted adjoint action is $\Ad_{R}(v) = Rv\tilde R$. By Theorem \claimref{std-algebra.spin-double-cover}, $\Ad$ maps the rotors onto $\mathrm{SO}(n)$, is a homomorphism with kernel $\{\pm 1\}$, so $\Ad_{R} = \Ad_{R'}$ exactly when $R' = \pm R$, and maps the versors of unit vectors onto $\mathrm{O}(n)$; it is continuous, being polynomial. \emph{The rotors form a compact set.} Let $\mathcal S$ be the set of even $R$ with $R\tilde R = 1$ such that $Rv\tilde R$ is a vector for every vector $v$. It is closed, and bounded because $\langle \tilde A A \rangle_{0}$ is a positive definite norm on $\Cl(n,0)$ (the sum of the squares of the blade coefficients), in which its elements have length one. Every rotor lies in $\mathcal S$. Conversely let $R \in \mathcal S$. In a finite-dimensional algebra $R\tilde R = 1$ gives $\tilde R R = 1$, so $(Rv\tilde R)^{2} = Rv^{2}\tilde R = Q(v)$: the map $v \mapsto Rv\tilde R$ is an isometry, and by the cited Theorem there is a versor $g$ of unit vectors with $\hat g\,v\,g^{-1} = Rv\tilde R$ for all $v$. Put $y = g^{-1}R$, so that $v\,y = \hat g^{-1}R\,v$. If $g$ is odd, $\hat g = -g$, $y$ is odd and $vy = -yv$ for every vector $v$; expanding $y$ in blades of an orthonormal basis, an odd blade commutes with each basis vector it contains, and conjugation by a basis vector maps each blade to itself up to sign, so no blade of $y$ survives and $y = 0$, impossible since $y$ is invertible (this is the kernel step of the proof of the cited Theorem). So $g$ is even, $y$ is even and commutes with every vector, hence is central, and by \claimref{primer.centre} the even part of the centre is $\R$ (the pseudoscalar, central for $n$ odd, is then odd): $y = \lambda$. Then $1 = R\tilde R = \lambda^{2}g\tilde g = \lambda^{2}$, $R = \pm g$, and $-g$ is again an even versor of unit vectors. So $\mathcal S$ is the set of rotors, which is therefore compact.

\emph{(L), the lifting.} Two consequences of the above. (a) $R$ and $-R$ are at distance $2$ in that norm. (b) For every $\epsilon > 0$ there is $\eta > 0$ such that, if $\Ad_{R}$ and $\Ad_{R'}$ are within $\eta$, then $R'$ is within $\epsilon$ of $R$ or of $-R$: otherwise there are sequences of rotors $R_{j}, R_{j}'$ with $\Ad_{R_{j}} - \Ad_{R_{j}'} \to 0$ and $R_{j}'$ at distance at least $\epsilon$ from $\pm R_{j}$; by compactness they converge along a subsequence to rotors $R, R'$ with $\Ad_{R} = \Ad_{R'}$ and $R' \neq \pm R$, a contradiction.

Take $\eta$ for $\epsilon = \tfrac12$ and, by uniform continuity, $\delta = 1/J$ such that $\kappa$ varies by less than $\eta$ on every square of side $\delta$. \emph{Paths.} Given a continuous path $A(t)$ in $\mathrm{SO}(n)$ varying by less than $\eta$ on intervals of length $\delta$, and a rotor $R_{0}$ over $A(0)$, define $R(t)$ on $[j\delta, (j+1)\delta]$ as the rotor over $A(t)$ within $\tfrac12$ of $R(j\delta)$; by (b) there is one and by (a) only one. It is continuous: for $t'$ near $t$ in the same interval, by (b) with a small $\epsilon'$ the rotor $R(t')$ is within $\epsilon'$ of $R(t)$ or of $-R(t)$, and the second is excluded because $R(t')$ and $R(t)$ are both within $\tfrac12$ of $R(j\delta)$, hence within $1$ of each other, while $-R(t)$ is at distance $2$ from $R(t)$. Two continuous lifts with the same start differ by a continuous sign, so the lift is unique. \emph{The square.} Lift $u \mapsto \kappa(u,0)$ from a rotor over $\kappa(0,0)$ to a path $\rho(u)$, and for each $u$ lift $t \mapsto \kappa(u,t)$ from $\rho(u)$ to $\hat\kappa(u,t)$. The map $\hat\kappa$ is continuous on the square: on the strip $j\delta \leq t \leq (j+1)\delta$, $\hat\kappa(u,t)$ is the rotor over $\kappa(u,t)$ within $\tfrac12$ of $\hat\kappa(u,j\delta)$, and if $\hat\kappa(\cdot, j\delta)$ is continuous then for $(u',t')$ near $(u,t)$ in the strip the rotor $\hat\kappa(u',t')$ is within $\epsilon'$ of $\pm\hat\kappa(u,t)$ and within $\tfrac12 + \tfrac12 + \epsilon'$ of $\hat\kappa(u,t)$, which excludes the minus sign; induction on $j$ starts from the continuity of $\rho$. Since $\kappa(u,1) = \kappa(u,0)$, $\hat\kappa(u,1) = \sigma(u)\,\hat\kappa(u,0)$ with $\sigma(u) = \pm 1$ continuous in $u$, hence constant. At $u = 1$ the path $\kappa(1,\cdot)$ is constant, its lift is constant by uniqueness, and $\sigma(1) = 1$: so $\hat\kappa(u,0) = \hat\kappa(u,1)$ for every $u$ and $\hat\kappa(1,t)$ is independent of $t$. A continuous path of rotors over $\kappa(0,\cdot)$ differs from $\hat\kappa(0,\cdot)$ by a continuous sign, so it is $\pm\hat\kappa(0,\cdot)$ and is closed. This proves (L).

\emph{(3).} Let $k$ be odd and suppose $E = G_{+} \oplus \underline{\R}^{N}$ is trivial over the sphere, $G_{+}$ the positive eigenbundle on $S^{(k)}$, of rank $2k$. $E$ carries the metric of $S^{(k)} \oplus \R^{N}$; a global frame, made orthonormal by the Gram-Schmidt process, which is continuous, gives a global orthonormal frame $\phi$. On the closed hemispheres $D_{N} = \{x_{3} \geq 0\}$ and $D_{S} = \{x_{3} \leq 0\}$ let $\nu$ and $\sigma$ be the orthonormal frames made of the normalised frames of (1), $k$ copies, and the standard basis of $\R^{N}$. Then $\phi = \nu\,\alpha_{N}$ on $D_{N}$ and $\phi = \sigma\,\alpha_{S}$ on $D_{S}$ with continuous $\alpha_{N}, \alpha_{S}$ into $\mathrm{O}(2k + N)$, and on the equator $\sigma = \nu\,\gamma$ with $\gamma(a)$ the turn of (2) in $k$ planes and the identity on the rest. So $\gamma = \alpha_{N}\alpha_{S}^{-1}$ on the equator. For $u, t \in [0,1]$ let $y_{N}(u,t)$ be the point of $D_{N}$ at angular distance $(1-u)\tfrac\pi2$ from the north pole and longitude $2\pi t$, $y_{S}(u,t)$ its mirror image in the equatorial plane, and
\[
c(u,t) = \alpha_{N}\bigl(y_{N}(u,t)\bigr)\,\alpha_{S}\bigl(y_{S}(u,t)\bigr)^{-1} .
\]
It is continuous, $c(u,0) = c(u,1)$, $c(1,t)$ is constant (the poles), and $c(0,t) = \gamma(2\pi t)$. Its determinant is continuous with values $\pm 1$ and equals $1$ at $c(0,0) = \gamma(0) = 1$, so $c$ maps into $\mathrm{SO}(2k + N)$. By (2) the rotor path $R(t)$, the product of the $k$ plane rotors at $a = 2\pi t$, is a continuous lift of $c(0,\cdot)$ with $R(0) = 1$ and $R(1) = (-1)^{k} = -1$. This contradicts (L), applied with $\kappa = c$ and $n = 2k + N$. So $G_{+} \oplus \underline{\R}^{N}$ is not trivial. The bundle $G_{-}$ is the pull-back of $G_{+}$ by the antipodal map $x \mapsto -x$, so a trivialisation of $G_{-} \oplus \underline{\R}^{N}$ would pull back to one of $G_{+} \oplus \underline{\R}^{N}$.

\emph{(4).} Let $h$ be a model completion standard outside the bounded set $U$, nondegenerate at every point of $\R^{3}$, with exterior form $K_{0}$ on $W_{0}$. By clause (4) of Theorem \claimref{synthesis.strict-completions-exist-exactly-when-the-signature-does-not-jump}, the base being contractible, $G_{+} \oplus (K_{0})_{+}$ is trivial on $\R^{3} \setminus U$. Choose $r$ with $U$ inside the ball of radius $r$. On the sphere of radius $r$ the positive eigenbundle of $g(x)$ is that of $g(x/r)$, the bundle $G_{+}$ of the unit sphere, and $(K_{0})_{+}$ is a fixed subspace, of dimension $N$ say. So $G_{+} \oplus \underline{\R}^{N}$ is trivial over the sphere, against (3). Strict completions: the signature of $G$ is $(2k,2k)$ at every point other than the origin, so clause (3) of that Theorem applies and the double is one; its $K = -G$ has $K_{\pm} = G_{\mp}$.

\emph{(5).} In $\Cl(4,0)$, $(e_{1}e_{2}e_{3}e_{4})^{2} = 1$, so $p_{4}$ is an idempotent, and $T = \Cl(4,0)p_{4}$ has dimension eight (computed). Left multiplications by $e_{1}, \dots, e_{4}$ on $T$ are symmetric for $\langle \tilde\psi\varphi \rangle_{0}$, as in the setting, square to one and anticommute. Right multiplication commutes with left multiplication, so $s \mapsto s\,p_{4}$ and $s \mapsto s\,e_{1}p_{4}$ are maps of $\Cl(3,0)$-modules from $S$ to $T$; each halves the product, and the images are orthogonal and together span $T$ (computed on the basis $p, e_{1}p, e_{2}p, e_{1}e_{2}p$). So $T$ is the orthogonal sum of two copies of $S$ with the product scaled by $\tfrac12$, a scaling that changes neither the symmetry of a map nor the eigenbundles of a family. On $T$, $(x + \beta e_{4})^{2} = Q(x) + \beta^{2}$, positive unless $x = 0$ and $\beta = 0$; with $\beta$ continuous, positive at the origin and zero outside $U$ (clause (8) of the cited Theorem), $x + \beta e_{4}$ is a continuous symmetric family, nondegenerate at every point, equal to $g(x)$ outside $U$: a model completion with $W_{0} = 0$. For $k$ even, $S^{(k)}$ is the orthogonal sum of $k/2$ copies of $T$ and the sum of the repairs is one. By clause (4) of the cited Theorem its eigenbundles, which are those of $G$ on the exterior, are trivial there.

\emph{(6).} (4) and (5) give the equivalence. Stable classes of bundles over the sphere add under direct sum; the class of $G_{+}$ on $S$ is not zero by (3) with $k = 1$, and its double is the class of $G_{+}$ on $S^{(2)}$, zero by (5): order exactly two. A repair on $S$ itself, of any form, is a model completion with $W_{0} = 0$ and is excluded by (4) with $k = 1$. A circle on the sphere bounds a disc of the sphere, over which $G_{+}$ is a bundle on a contractible space, so its restriction to the circle is trivial: the obstruction is not a class of degree one. Separately, the linear system for a symmetric map on $S$ anticommuting with $e_{1}, e_{2}, e_{3}$ has only the zero solution (computed, as in clause (8) of the cited Theorem for $m = 3$), which excludes the repairs of the special form $x + \beta E$ and nothing more.
